% Supplement, Part II, section S8: vanishing and rate of the remainder.
% Sources and corrections: see CONC_s8s9.md.
\section{Vanishing and rate of the remainder}\label{sec:remainder}

This section and the next prove the first links of Theorem~2.3 of the main text. Theorem~\ref{thm:remainder-vanish} identifies the constant term of the bounded remainder of Theorem~R (\S\ref{sec:rate}): it is the equal-height mass $B_{d,L}$ for the candidate density and $0$ for the elder density. Theorem~\ref{thm:remainder-rate} shows that the elder remainder is $O(\ell^{1/4})$. Both theorems hold for every fixed dimension $d\ge2$ and side $L>0$; at $d=3$, $L=24$ the elder density is $\nu_{3,24}$.

\subsection{Setting and inputs}\label{ssec:remainder-setting}

\begin{convention}[notation for Part~II]\label{conv:remainder-notation}
\begin{enumerate}[(a),leftmargin=2em,itemsep=2pt]
\item \emph{Field and constants.} Fix $d\ge2$ and $L>0$, and let $f$ be the centred stationary Gaussian field of variance one on the torus $\T^d_L=\R^d/(L\Z^d)$ with covariance $K_L(z)=\sum_{n\in\Z^d}e^{-|z+Ln|^2/2}\big/\sum_{n\in\Z^d}e^{-|Ln|^2/2}$; for $d=3$, $L=24$ this is (1.1) of the main text. Put $m:=d-1$. Unless stated otherwise, $C,c,N$ are positive constants depending on $d$ and $L$ only; they may change from line to line. The index of a symmetric matrix is its number of negative eigenvalues.
\item \emph{Densities.} $\nu_{\cand}(\ell)$ is the $\ell$-density, per unit volume, of the expected number of ordered pairs $(M,S)$ of a local maximum $M$ and a critical point $S$ of index $d-1$ with $f(M)-f(S)=\ell>0$, all births, gaps and separations included; $\nu_{\eld}(\ell)$ is the same density for the pairs that are finite ordinary superlevel $H_0$ bars (the essential class excluded); $\rho_{\rej}:=\nu_{\cand}-\nu_{\eld}$. All three are taken in their canonical marked Kac--Rice density versions (Convention~2.2 of the main text and \S\ref{sec:setting}; for general $d,L$ the same construction, \cite[\S9]{parentP}; Proposition~\ref{prop:parent-KR}(b), \eqref{eq:parent-135} and Lemma~\ref{lem:parent-far}). At $d=3$, $L=24$, $\nu_{\eld}=\nu_{3,24}$. A superscript $\mathrm{near},\rho$, $\mathrm{far},\rho$ or $[\rho,\rho')$ restricts to pairs with $\dist(M,S)<\rho$, $\ge\rho$ or in $[\rho,\rho')$. The constant $c_{d,L}>0$ is the leading coefficient of \cite[(15.2)]{parentP} (Proposition~\ref{prop:parent-coef}), so that $\nu_{\cand}(\ell)\sim\nu_{\eld}(\ell)\sim c_{d,L}\ell^{-1/3}$ \cite[Theorem~C]{parentP} (Theorem~\ref{thm:parent-C}); at $d=3$, $L=24$ it equals $c_{3,24}$ of (2.9) of the main text, both being the leading coefficient of $\nu_{3,24}$.
\item \emph{Pinned pairs.} By stationarity the midpoint of the pair is the origin. For a direction $u\in S^{d-1}$ (the $t$ of (3.4) of the main text; we write $u$, as in (2.13) of the main text), an orthonormal frame $\Theta$ of $u^\perp$, and coordinates $(x,y)\in\R\times\R^m$ along $u$ and $\Theta$, the pins are
\[
M=-\tfrac r2u,\quad S=\tfrac r2u,\qquad f(M)=b,\quad f(S)=b-kr^3,\qquad \nabla f(M)=\nabla f(S)=0,
\]
with $b\in\R$ and gap mark $k>0$; in the notation of (3.5) of the main text $k=\kappa/6$, and the lifetime is $\ell=kr^3$. Put $P:=1+|b|+k$. The observation vector $U_r$ has the axial rows $\tfrac12(f(M)+f(S))$, $(f(S)-f(M))/r$, $(\partial_xf(S)-\partial_xf(M))/r$, $\tfrac6{r^2}[\partial_xf(M)+\partial_xf(S)-2(f(S)-f(M))/r]$ and the transverse rows $\tfrac12(\partial_{y_j}f(M)+\partial_{y_j}f(S))$, $(\partial_{y_j}f(S)-\partial_{y_j}f(M))/r$; the map from the raw pins $(f(M),\nabla f(M),f(S),\nabla f(S))$ to $U_r$ has absolute determinant $12r^{-(d+3)}$, and the target of $U_r$ is $v_r=(b-kr^3/2,\,-kr^2,\,0,\,12k,\,0,\dots,0)$. At $r=0$, $U_0=(f,\partial_xf,\partial_x^2f,\partial_x^3f,\nabla_yf,\nabla_y\partial_xf)(0)$ with target $v_0=v_0(b,k)=(b,0,0,12k,0,\dots,0)$. We write $\pi_r$, $\pi_0$ for the Gaussian densities of $U_r$, $U_0$, $Q=Q_{r,b,k}$ for the continuous Gaussian regression law of the field given $U_r=v_r$, and call the regression law given $U_0=v_0$ the \emph{contact law}.
\item \emph{Weights.} For a symmetric matrix $H$ let $F_j(H):=|\det H|$ if $H$ has exactly $j$ negative eigenvalues and $F_j(H):=0$ otherwise (so $F_j(H)=0$ if $H$ is singular). With $H_i=D^2f(i)$ ($i\in\{M,S\}$), the typed weight is $W_r:=|\det H_M\det H_S|\,\ind\{H_M\prec0,\ \Ind H_S=d-1\}$ and $Z_r:=\E_QW_r$. The scaled endpoint Hessians are
\[
K_i:=D_r^{-1}H_iD_r^{-1}=\begin{pmatrix}\alpha_i&\sqrt r\,\beta_i^{\top}\\ \sqrt r\,\beta_i&A_i\end{pmatrix},\qquad \alpha_i=\frac{\partial_x^2f(i)}r,\quad \beta_i=\frac{\nabla_y\partial_xf(i)}r,\quad A_i=D_y^2f(i),
\]
with $D_r=\operatorname{diag}(\sqrt r,I_m)$; thus $|\det H_i|=r|\det K_i|$, $K_i$ is congruent to $H_i$, and $W_r/r^2=F_d(K_M)F_{d-1}(K_S)$.
\item \emph{Elder mark.} For $M\in\T^d_L$, $d_f(M):=\sup\{\min_{s\in[0,1]}f(\gamma(s)):\gamma\in C([0,1];\T^d_L),\ \gamma(0)=M,\ f(\gamma(1))>f(M)\}$ (with $\sup\emptyset=-\infty$) is the maximin connection level, and $e:=\ind\{d_f(M)=f(S)\}$. The mark $e$ is Borel \cite[\S8]{parentP} (Lemma~\ref{lem:parent-mark}; cf.\ \S\ref{sec:selection}); on the Morse locus with distinct critical values, $e=1$ exactly when $S$ is the elder death saddle of $M$.
\item \emph{Kernels.} $\mathcal A_r(b,k,u):=12\pi_r(v_r)\E_Q[W_r/r^2]$, $\mathcal A^{\eld}_r:=12\pi_r(v_r)\E_Q[(W_r/r^2)e]$ and $\mathcal A^{\rej}_r:=\mathcal A_r-\mathcal A^{\eld}_r=12\pi_r(v_r)\E_Q[(W_r/r^2)(1-e)]$. At contact, $A_0$ denotes the transverse Hessian $D_y^2f(0)$ under the contact law, $z_0(b,k,u):=\E[(6k)^2(\det A_0)^2\ind\{A_0\prec0\}]$ and $\mathcal A_0(b,k,u):=12\pi_0(v_0)z_0$. Finally $H(b,k):=CP^Ne^{-c(b^2+k^2)}$, with $N$ enlarged when necessary.
\item \emph{Two-site kernels.} For $y\in\T^d_L\setminus\{0\}$ let $O_y=(f(0),\nabla f(0),f(y),\nabla f(y))$, $v_{b,\ell}=(b,0,b-\ell,0)$, $p_y$ the density of $O_y$, $Q_{y,b,\ell}$ the regression law given $O_y=v_{b,\ell}$, $W=|\det H_0\det H_y|\,\ind\{H_0\prec0,\ \Ind H_y=d-1\}$ and $e_y=\ind\{d_f(0)=f(y)\}$. The unmarked candidate kernel, the equal-height kernel and the equal-height mass are
\[
\Psi^{\cand}_\ell(b,y)=p_y(v_{b,\ell})\,\E_{Q_{y,b,\ell}}[W],\qquad \Psi_0(b,y)=p_y(v_{b,0})\,\E_{Q_{y,b,0}}[W],
\]
\[
B_{d,L}=\int_{\T^d_L\setminus\{0\}}\int_\R\Psi_0(b,y)\,db\,dy .
\]
\item \emph{Radii.} $r_R\in(0,\frac12]$ is the radius of Proposition~\ref{prop:remainder-R}, $r_K$ that of Proposition~\ref{prop:remainder-K}, and $r_Z:=r_K$. The \emph{admissible radius} is $r_0^*:=\min\{r_R,r_K,r_Z\}$, decreased if necessary so that $r_0^*<L/(4\sqrt2)$, the embedding radius of the cap cylinder (each of the two radius statements remains true for a smaller radius; for $L=24$ nothing changes, as $r_R\le1<24/(4\sqrt2)$). The dominating variable is
\[
T:=1+k+\|F_r\|_{C^7}+\|F_0\|_{C^7}+r^{-2}\|F_r-F_0\|_{C^2},
\]
for the coupling $(F_r,F_0)$ of Proposition~\ref{prop:remainder-R}.
\end{enumerate}
\end{convention}

The marked Kac--Rice disintegration (\S\ref{sec:contact}; \cite[\S\S9--10, (10.2)]{parentP}; Propositions~\ref{prop:parent-KR} and~\ref{prop:parent-ledger}) gives the candidate intensity $r\mathcal A_r\,dr\,db\,dk\,d\sigma(u)$ per unit volume, with $\sigma$ the surface measure, and the elder and rejected intensities with $\mathcal A^{\eld}_r$, $\mathcal A^{\rej}_r$ in place of $\mathcal A_r$. Pushing forward at fixed $r$ with $dk=d\ell/r^3$ gives, for $0\le\rho<\rho'\le r_R$, $\ell>0$ and $\bullet\in\{\cand,\eld,\rej\}$ (with $\mathcal A^{\cand}_r:=\mathcal A_r$),
\begin{equation}\label{eq:remainder-sep}
\nu_\bullet^{[\rho,\rho')}(\ell)=\int_\rho^{\rho'}\!\!\int_\R\int_{S^{d-1}}r^{-2}\mathcal A^\bullet_r(b,\ell/r^3,u)\,d\sigma(u)\,db\,dr .
\end{equation}
The leading coefficient is $c_{d,L}=\int_{S^{d-1}}\int_\R\int_0^\infty\mathcal A_0(b,k,u)\,(3k^{2/3})^{-1}\,dk\,db\,d\sigma(u)$ (\S\ref{sec:rate}; \cite[(13.6)]{parentP}; Proposition~\ref{prop:parent-near}); the substitution $k=\ell/r^3$, with $|dk/dr|=3\ell r^{-4}$ and the order of the endpoints reversed, turns it into
\begin{equation}\label{eq:remainder-lead}
c_{d,L}\,\ell^{-1/3}=\int_0^\infty\!\!\int_\R\int_{S^{d-1}}r^{-2}\mathcal A_0(b,\ell/r^3,u)\,d\sigma(u)\,db\,dr .
\end{equation}
Translating the pins to $0$ and $y=ru$, the polar factor $r^{d-1}$, the pin density $12r^{-(d+3)}\pi_r(v_r)$ and the determinant scale $r^2$ give, for $0<r\le r_0^*$,
\begin{equation}\label{eq:remainder-polar}
r^{-2}\mathcal A_r(b,\ell/r^3,u)=r^{d-1}\Psi^{\cand}_\ell(b,ru),\qquad r^{-2}\mathcal A_r(b,0,u)=r^{d-1}\Psi_0(b,ru),
\end{equation}
where $\mathcal A_r(b,0,u)$ is defined with the pinned law at $k=0$. On $\{\dist(0,y)\ge\rho\}$ the height map $(b,\ell)\mapsto(b,b-\ell)$ has absolute Jacobian one, and
\begin{equation}\label{eq:remainder-far}
\nu_{\cand}^{\mathrm{far},\rho}(\ell)=\int_{\dist(0,y)\ge\rho}\int_\R\Psi^{\cand}_\ell(b,y)\,db\,dy,\qquad
\nu_{\eld}^{\mathrm{far},\rho}(\ell)=\int_{\dist(0,y)\ge\rho}\int_\R p_y(v_{b,\ell})\,\E_{Q_{y,b,\ell}}[We_y]\,db\,dy .
\end{equation}
Separation boundaries are Lebesgue-null in the pair variables, so the ranges in \eqref{eq:remainder-sep}--\eqref{eq:remainder-far} may be combined freely.

We use three groups of results that are proved elsewhere.

\begin{proposition}[the contact coupling and the bounds of Theorem~R]\label{prop:remainder-R}
There is $r_R\in(0,\frac12]$, with $r_R\le r_0$ of \eqref{eq:rate-r0}, such that, for $0<r\le r_R$, $b\in\R$ and every frame:
\begin{enumerate}[(i),leftmargin=2em]
\item Let $F$ be an unconditioned copy of the field, $\Sigma_r=\Cov(U_r)$ and $C_r(z)=\Cov(F(z),U_r)$. Then $F_r:=F+C_r\Sigma_r^{-1}(v_r-U_r(F))$ has law $Q$ and $F_0:=F+C_0\Sigma_0^{-1}(v_0-U_0(F))$ has the contact law. The matrices $\Sigma_r$ are uniformly positive definite and bounded, $\Sigma_r-\Sigma_0$, $\Sigma_r^{-1}-\Sigma_0^{-1}$ and $\|C_r-C_0\|_{C^q}$ are $O(r^2)$, and these covariance facts do not depend on the target, so (i) and (ii) hold for all $k\ge0$. For every finite $p\ge1$ and every derivative order $q$,
\begin{equation}\label{eq:remainder-R4}
\|F_r-F_0\|_{L^p(C^q)}\le C_{p,q}\,r^2P,\qquad \bigl\|1+\|F_r\|_{C^q}+\|F_0\|_{C^q}\bigr\|_{L^p}\le C_{p,q}\,P .
\end{equation}
In particular $\|T\|_{L^p}\le C_pP$ for every $p$.
\item $\pi_r(v_r)\le Ce^{-c(b^2+k^2)}$ and $|\pi_r(v_r)-\pi_0(v_0)|\le Cr^2P^2e^{-c(b^2+k^2)}$.
\item For a symmetric $m\times m$ matrix $A$, $\alpha\in\R$, $\beta\in\R^m$ and $t\ge0$, the matrix $K_t=\bigl(\begin{smallmatrix}\alpha&\sqrt t\beta^{\top}\\ \sqrt t\beta&A\end{smallmatrix}\bigr)$ satisfies, also for singular $A$,
\begin{equation}\label{eq:remainder-block}
\det K_t=\alpha\det A-t\,\beta^{\top}\adj(A)\,\beta .
\end{equation}
\item For $k\ge0$, with $f=F_r$ and $A_0:=D_y^2F_0(0)$: $\|A_i-A_0\|\le(\frac r2+r^2)T\le rT$, $\|A_i\|+\|A_0\|+\|\beta_i\|\le CT$, $|\alpha_M+6k|+|\alpha_S-6k|\le CrT$, $F_d(K_M),F_{d-1}(K_S)\le C(k+r)T^d$, and
\begin{equation}\label{eq:remainder-R11}
\frac{Z_r}{r^2}\le C(k+r)^2P^N,\qquad 0\le\mathcal A_r\le(k+r)^2H,\qquad 0\le\mathcal A_0\le k^2H,\qquad |\mathcal A_r-\mathcal A_0|\le r(k+r)H .
\end{equation}
\end{enumerate}
\end{proposition}

\begin{proof}
These are the coupling, the pin-density estimates, the block-determinant lemma and the weight bounds in the proof of Theorem~R (\S\ref{sec:rate}); Theorem~R and its proof hold for every fixed $d\ge2$ and $L>0$. These bounds hold also at $k=0$: no division by $k$ or by $Z_r$ occurs in the proofs of \eqref{eq:rate-R10}--\eqref{eq:rate-R11} (see the remark after \eqref{eq:rate-R11}); the covariance $\Sigma_r$ does not depend on the target, so the Gaussian bounds hold with $k=0$; the endpoint identities \eqref{eq:parent-52} of Lemma~\ref{lem:parent-endpoint} are deterministic and hold for every value of the gap mark, including $k=0$; and the block-determinant comparison (iii) holds also for singular blocks. The variable $T$ dominates the dominating variable used there (which is built from $C^4$ norms), so its bounds hold with $T$ after adjusting $C$. By Taylor's formula, $\|A_i-A_0\|\le\frac r2\|F_r\|_{C^3}+\|F_r-F_0\|_{C^2}\le(\frac r2+r^2)T$, and $\frac r2+r^2\le r$ as $r\le\frac12$; a smaller radius keeps every statement.
\end{proof}

\begin{proposition}[gap-mark-explicit weight bounds]\label{prop:remainder-K}
Let $G_r:=\{\lambda_{\min}(-A_M)>4rM_3^2/(3k)\}\cap\{rM_4\le3k/10\}$ be the cap region, where $M_3,M_4$ are the local $C^3$ and $C^4$ seminorms of the field on the cap cylinder \cite[\S\S4,7]{parentP} (Definition~\ref{def:cap-norms}). There are $r_K>0$, $C$ and $N$ such that, for every $b\in\R$, $k>0$, every frame and $0<r\le r_K$,
\begin{equation}\label{eq:remainder-K}
\E_Q[W_r/r^2]\le C(k+r)^2P^N,\qquad \E_Q\bigl[(W_r/r^2)\ind_{G_r^c}\bigr]\le C\,\frac{r^3}k\,P^N\quad\text{whenever }r\le\min\{k,r_K\}.
\end{equation}
Moreover, if also $r<L/(4\sqrt2)$, then $1-e\le\ind_{G_r^c}$ $Q$-almost surely on $\{W_r>0\}$.
\end{proposition}

\begin{proof}
The first bound is that of \eqref{eq:remainder-R11}, as $Z_r=\E_QW_r$, and the second is Lemma~\ref{lem:remainder-capK}; take $r_K:=\min\{r_R,r_0\}$ with the radius $r_0$ of \eqref{eq:rate-r0}. Together the two bounds are Theorem~K of record [R7] of the main text. The last statement is the deterministic cap theorem \cite[\S\S7--8]{parentP} (Lemma~\ref{lem:parent-capimpl} and Theorem~\ref{thm:cap-main}), which needs the cap cylinder to be embedded, that is, $r<L/(4\sqrt2)$ (Definition~\ref{def:parent-pinned}): on $G_r$ and on the typed support, the global elder partner of $M$ is $S$, and $Q$ is almost surely Morse with distinct critical values \cite[\S8]{parentP} (Lemma~\ref{lem:parent-generic}(a)).
\end{proof}

\begin{proposition}[the equal-height kernel]\label{prop:remainder-Z}
The following hold.
\begin{enumerate}[(i),leftmargin=2em]
\item $0<B_{d,L}<\infty$, and, for the canonical marked Kac--Rice density versions, $\rho_{\rej}(\ell)\to B_{d,L}$ as $\ell\downarrow0$.
\item For fixed $y\ne0$ and $b$, the covariance of $Q_{y,b,\ell}$ does not depend on $\ell$; on one probability space $f_\ell=f_0+\ell h_y$, with $f_0$ of law $Q_{y,b,0}$ and $h_y$ a deterministic smooth function with $h_y(0)=0$, $h_y(y)=-1$, $\nabla h_y(0)=\nabla h_y(y)=0$. The two endpoint Hessians are almost surely nonsingular, and $\Psi^{\cand}_\ell(b,y)\to\Psi_0(b,y)$ as $\ell\downarrow0$ for every $y\ne0$ and $b\in\R$.
\item For every $\rho\in(0,L/2]$ there are $C_\rho,c_\rho>0$ with $0\le\Psi^{\cand}_\ell(b,y)\le C_\rho(1+|b|)^{2d}e^{-c_\rho b^2}$ for $\dist(0,y)\ge\rho$ and $0<\ell\le1$.
\end{enumerate}
\end{proposition}

\begin{proof}
Item (iii) is the nondegenerate far bound \cite[\S14, (14.1)]{parentP} (Lemma~\ref{lem:parent-far}). Items (i) and (ii) are Theorem~Z of record [R7] of the main text; they are proved in \S\ref{ssec:remainder-Z}, item (ii) and the bounds $0<B_{d,L}<\infty$ directly, and the limit in (i) from \eqref{eq:remainder-Rp1} and Theorem~\ref{thm:remainder-rate}, whose proofs use this proposition only through (ii), (iii) and the finiteness of $B_{d,L}$ (the proof of \eqref{eq:remainder-Rp2} then uses the limit).
\end{proof}

\subsection{The even contact expansion}\label{ssec:remainder-even}

\begin{lemma}[even contact expansion]\label{lem:remainder-even}
There are $C,N$ such that for all $b\in\R$, $k>0$, $u\in S^{d-1}$ and $0<r\le r_0^*$,
\begin{equation}\label{eq:remainder-E}
\Bigl|\frac{Z_r}{r^2}-z_0(b,k,u)\Bigr|\le Cr^2P^N,\qquad |\mathcal A_r-\mathcal A_0|\le r^2H(b,k).
\end{equation}
\end{lemma}

The first-order term of \eqref{eq:remainder-R11} is absent: the two endpoint corrections enter the typed product with opposite signs and cancel identically (Steps~1--3 below), and on the layers where the inertia can flip the product of the two typed determinants is $O(r^2)$ (Step~4).

\begin{proof}
Work under the coupling of Proposition~\ref{prop:remainder-R}(i). Write $f=F_r$, $f_j=\partial_x^jf(0)$, and $O(x)$ for a quantity bounded by $C|x|$ with $C=C(d)$; $\|T\|_{L^p}\le C_pP$ by \eqref{eq:remainder-R4} with $q=7$ and $q=2$.

\emph{Step 1 (axial pins to second order).} The third and fourth rows of $U_r$ take the pinned values $0$ and $12k$. Taylor's formula with integral remainder (the rules are centred, so the odd orders cancel) gives
\[
\begin{gathered}
\frac{\partial_xf(S)-\partial_xf(M)}r=f_2+\frac{r^2}{24}f_4+O(r^4T)=0,\\
\frac6{r^2}\Bigl[\partial_xf(M)+\partial_xf(S)-\frac{2(f(S)-f(M))}r\Bigr]=f_3+\frac{r^2}{40}f_5+O(r^4T)=12k .
\end{gathered}
\]
Substituting both identities in $\partial_x^2f(\mp r/2)=f_2\mp\tfrac r2f_3+\tfrac{r^2}8f_4+O(r^3T)$ gives
\begin{equation}\label{eq:remainder-alpha}
\alpha_M=-6k+\frac r{12}f_4+\rho_M,\qquad \alpha_S=6k+\frac r{12}f_4+\rho_S,\qquad |\rho_M|+|\rho_S|\le Cr^2T .
\end{equation}
(For a quintic profile the expansion is exact: $\alpha_{M,S}=\mp6k+\frac r{12}f_4\mp\frac{r^2}{120}f_5$.) The first-order coefficients coincide, so
\begin{equation}\label{eq:remainder-alphaprod}
\begin{aligned}
\alpha_M\alpha_S&=-36k^2+6k(\rho_M-\rho_S)+r^2\Bigl(\frac{f_4}{12}\Bigr)^2+\frac{rf_4}{12}(\rho_M+\rho_S)+\rho_M\rho_S\\
&=-36k^2+O\bigl(r^2(1+k)T^2\bigr).
\end{aligned}
\end{equation}

\emph{Step 2 (transverse blocks).} Put $B:=\partial_xD_y^2f(0)$. Since $\|D_y^2f(0)-A_0\|\le r^2T$ by the coupling,
\begin{equation}\label{eq:remainder-A}
A_M=A_0-\tfrac r2B+O(r^2T),\qquad A_S=A_0+\tfrac r2B+O(r^2T),
\end{equation}
and $\det(A_0+V)=\det A_0+\tr(\adj(A_0)V)+O(\|V\|^2T^{m-2})$ gives $\det A_M\det A_S=(\det A_0)^2+\det A_0\,[-\tfrac r2+\tfrac r2]\tr(\adj(A_0)B)+O(r^2T^{2m})=(\det A_0)^2+O(r^2T^{2m})$. The transverse difference row gives $(\nabla_yf(S)-\nabla_yf(M))/r=\nabla_y\partial_xf(0)+\frac{r^2}{24}\nabla_y\partial_x^3f(0)+O(r^4T)=0$, so $\nabla_y\partial_xf(0)=O(r^2T)$; with $\gamma:=\nabla_y\partial_x^2f(0)$,
\begin{equation}\label{eq:remainder-beta}
\begin{gathered}
\beta_M=-\tfrac12\gamma+O(rT),\qquad \beta_S=\tfrac12\gamma+O(rT),\\
q_i:=\beta_i^{\top}\adj(A_i)\beta_i=q+O(rT^{m+1}),\qquad q:=\tfrac14\gamma^{\top}\adj(A_0)\gamma .
\end{gathered}
\end{equation}
The sign of $\beta_i$ enters $q_i$ quadratically. That $\nabla_y\partial_xf(0)$ is $O(r^2)$ and not $O(r)$ is essential: an $O(r)$ value would shift $\beta_M$ and $\beta_S$ by a common $O(1)$ vector and make $q_M-q_S=O(1)$.

\emph{Step 3 (the product of determinants).} By \eqref{eq:remainder-block} with $t=r$ (the off-diagonal entries of $K_i$ are $\sqrt r\beta_i$), $\det K_i=\alpha_i\det A_i-rq_i$ exactly. Hence
\[
\det K_M\det K_S=\alpha_M\alpha_S\det A_M\det A_S-r\bigl(\alpha_M\det A_M\,q_S+\alpha_S\det A_S\,q_M\bigr)+r^2q_Mq_S .
\]
By \eqref{eq:remainder-alphaprod} and Step~2 the first term is $-36k^2(\det A_0)^2+O(r^2(1+k)^2T^{N_1})$. By \eqref{eq:remainder-alpha}, \eqref{eq:remainder-A} and \eqref{eq:remainder-beta}, $\alpha_M\det A_M\,q_S=-6k\det A_0\,q+O(r(1+k)T^{N_1})$ and $\alpha_S\det A_S\,q_M=+6k\det A_0\,q+O(r(1+k)T^{N_1})$, so the middle term is $O(r^2(1+k)T^{N_1})$; the last term is $O(r^2T^{N_1})$. Thus
\begin{equation}\label{eq:remainder-prod}
\bigl|\det K_M\det K_S+36k^2(\det A_0)^2\bigr|\le Cr^2(1+k)^2T^{N_1}.
\end{equation}
Structurally, $\det K_M=-6k\det A_0+rY+O(r^2)$ and $\det K_S=+6k\det A_0+rY+O(r^2)$ with the \emph{same}
\begin{equation}\label{eq:remainder-Y}
Y:=\frac{f_4}{12}\det A_0+3k\,\tr\bigl(\adj(A_0)B\bigr)-q,
\end{equation}
so the $r$-term of the product is $6k\det A_0\,(Y-Y)=0$. This is the reflection symmetry of the pinned family: reversing the axis exchanges $M$ and $S$, and the rows of $U_r$ are even in $r$.

\emph{Step 4 (inertia).} For invertible $A_i$, Haynsworth's inertia additivity gives $\operatorname{inertia}(K_i)=\operatorname{inertia}(A_i)+\operatorname{inertia}(\sigma_i)$ with $\sigma_i:=\alpha_i-r\beta_i^{\top}A_i^{-1}\beta_i$. So $\Ind K_M=d$ if and only if $A_M\prec0$ and $\sigma_M<0$, and $\Ind K_S=d-1$ if and only if either $A_S\prec0$ and $\sigma_S>0$, or $\Ind A_S=m-1$ and $\sigma_S<0$. Put $\mathsf u:=C_1rT(1+T/k)$ with $C_1=C_1(d)\ge2$ large, and let $G:=\{\lambda_{\max}(A_0)\le-\mathsf u\}\cap\{k\ge C_1rT\}$. On $G$: $\|A_i-A_0\|\le rT\le\mathsf u/C_1$, so $A_M,A_S\prec0$; $\alpha_S\ge6k-CrT>0$ and $r\beta_S^{\top}(-A_S)^{-1}\beta_S\ge0$, so $\sigma_S>0$; $\alpha_M\le-6k+CrT\le-5k$ and $r\beta_M^{\top}(-A_M)^{-1}\beta_M\le rT^2/\lambda_{\min}(-A_M)\le2rT^2/\mathsf u\le2k/C_1$, so $\sigma_M<0$. Hence on $G$ the indices are $d$ and $d-1$, $\ind\{A_0\prec0\}=1$, and
\[
\bigl|F_d(K_M)F_{d-1}(K_S)-36k^2(\det A_0)^2\ind\{A_0\prec0\}\bigr|=\bigl||\det K_M\det K_S|-36k^2(\det A_0)^2\bigr|
\]
is at most the left side of \eqref{eq:remainder-prod}. Off $G$ there are three cases, in each of which both typed products are $r^2$ times a polynomial in $T$:
\begin{itemize}[leftmargin=1.5em]
\item on $\{k<C_1rT\}$: $F_d(K_M)F_{d-1}(K_S)\le C(k+r)^2T^{2d}\le Cr^2T^{2d+2}$ by Proposition~\ref{prop:remainder-R}(iv), and $36k^2(\det A_0)^2\le Cr^2T^{2m+2}$;
\item on $\{\lambda_{\max}(A_0)\ge\mathsf u\}\cap\{k\ge C_1rT\}$: $\lambda_{\max}(A_M)\ge\mathsf u-rT>0$, so $K_M$ is not negative definite, $F_d(K_M)=0$, and $\ind\{A_0\prec0\}=0$; both products vanish;
\item on $\{|\lambda_{\max}(A_0)|<\mathsf u\}\cap\{k\ge C_1rT\}$: by Weyl's inequality $|\det A_i|\le(|\lambda_{\max}(A_0)|+rT)T^{m-1}\le2\mathsf uT^{m-1}$, so $|\det K_i|\le(6k+CrT)\,2\mathsf uT^{m-1}+rT^{m+1}\le Ck\mathsf uT^{m-1}$ (as $rT^2\le k\mathsf u/C_1$); hence $F_d(K_M)F_{d-1}(K_S)\le Ck^2\mathsf u^2T^{2m-2}$ and likewise $36k^2(\det A_0)^2\ind\{A_0\prec0\}\le36k^2\mathsf u^2T^{2m-2}$; with $k^2\mathsf u^2=C_1^2r^2T^2(k+T)^2$ both are $\le Cr^2(k+T)^2T^{2m}$.
\end{itemize}
No smallness of the probability of these layers is needed: on the last layer each typed determinant is $O(r)$.

\emph{Step 5 (assembly).} Taking expectations, $|Z_r/r^2-z_0|\le\E[Cr^2(1+k)^2T^{N_1}+Cr^2(k+T)^2T^{2m+2}]\le Cr^2P^N$. For $\mathcal A_r-\mathcal A_0=12[\pi_r(v_r)-\pi_0(v_0)]Z_r/r^2+12\pi_0(v_0)[Z_r/r^2-z_0]$ use Proposition~\ref{prop:remainder-R}(ii), the bound $Z_r/r^2\le C(k+r)^2P^N$ of \eqref{eq:remainder-R11} and the first bound: every term is $r^2$ times a polynomial in $P$ times $e^{-c(b^2+k^2)}$, that is, at most $r^2H$.
\end{proof}

\subsection{The constant term of the remainder}\label{ssec:remainder-vanish}

\begin{theorem}[Theorem R$^+$: the constant term of the remainder]\label{thm:remainder-vanish}
For every fixed $d\ge2$ and $L>0$, and for the canonical marked Kac--Rice density versions (the versions for which Proposition~\ref{prop:remainder-Z}(i) is stated; $\nu_{\cand}$ and $\nu_{\eld}$ are the unrestricted densities), as $\ell\downarrow0$,
\begin{align}
\nu_{\cand}(\ell)&=c_{d,L}\,\ell^{-1/3}+B_{d,L}+o(1),\label{eq:remainder-Rp1}\\
\nu_{\eld}(\ell)&=c_{d,L}\,\ell^{-1/3}+o(1)=c_{d,L}\,\ell^{-1/3}\bigl(1+o(\ell^{1/3})\bigr).\label{eq:remainder-Rp2}
\end{align}
Moreover, for every $r_0\in(0,r_0^*]$, as $\ell\downarrow0$,
\begin{equation}\label{eq:remainder-Rp3}
\begin{gathered}
\nu_{\cand}^{\mathrm{near},r_0}(\ell)-c_{d,L}\,\ell^{-1/3}\to\int_{0<\dist(0,y)<r_0}\int_\R\Psi_0(b,y)\,db\,dy,\\
\bigl|\nu_{\cand}^{\mathrm{near},r_0}(\ell)-c_{d,L}\,\ell^{-1/3}\bigr|\le C\bigl(r_0+\ell^2r_0^{-7}\bigr),
\end{gathered}
\end{equation}
the bound holding for every $\ell>0$.
\end{theorem}

So the bounded remainder of Theorem~R has the constant term $B_{d,L}$ for the candidate density and the constant term $0$ for the elder density: no expansion $\nu_{\eld}=c_{d,L}\ell^{-1/3}+C_0+o(1)$ with $C_0\ne0$ is asserted, and the coefficient $c_1$ of \S\ref{sec:cusp} is the coefficient of $\ell^{1/4}$, a different term. The theorem gives no rate in \eqref{eq:remainder-Rp1}, no numerical constants, no uniformity in $d$ or $L$, and no statement about other almost-everywhere-equal versions of the densities. It does not enlarge the scope of Proposition~\ref{prop:remainder-Z}(i); it identifies the constant term of the remainder of Theorem~R by means of it. A rate for \eqref{eq:remainder-Rp2} is Theorem~\ref{thm:remainder-rate}.

\begin{proof}
\emph{The separation variable.} By \eqref{eq:remainder-sep} and \eqref{eq:remainder-lead},
\begin{equation}\label{eq:remainder-diff}
\begin{aligned}
\nu_{\cand}^{\mathrm{near},r_0}(\ell)-c_{d,L}\ell^{-1/3}&=\int_0^{r_0}\!\!\int\!\!\int r^{-2}(\mathcal A_r-\mathcal A_0)(b,\ell/r^3,u)\,d\sigma\,db\,dr\\
&\quad-\int_{r_0}^\infty\!\!\int\!\!\int r^{-2}\mathcal A_0(b,\ell/r^3,u)\,d\sigma\,db\,dr .
\end{aligned}
\end{equation}

\emph{Domination.} By Lemma~\ref{lem:remainder-even}, $|r^{-2}(\mathcal A_r-\mathcal A_0)(b,\ell/r^3,u)|\le H(b,\ell/r^3)\le Ce^{-cb^2}$, an integrable majorant on $(0,r_0)\times\R\times S^{d-1}$ that does not depend on $\ell$. (With the first-order bound of \eqref{eq:remainder-R11} the integrand is only bounded by $(1+\ell/r^4)H$.) The second term of \eqref{eq:remainder-diff} is at most $\int_{r_0}^\infty r^{-2}(\ell/r^3)^2H\,dr\le C\ell^2r_0^{-7}$, because $\mathcal A_0\le k^2H$.

\emph{Pointwise limit.} Fix $r,b,u$. As $\ell\downarrow0$, $k=\ell/r^3\downarrow0$ and $\mathcal A_r(b,k,u)\to\mathcal A_r(b,0,u)$: the pinned law is the regression law with target $v_r(b,k)$, affine in the target with a $k$-free covariance, so $F_r$ depends continuously on $k$ in $C^2$; the typed weight $W_r=F_d(H_M)F_{d-1}(H_S)$ is continuous in the endpoint Hessians, as $F_d$ and $F_{d-1}$ are locally Lipschitz (Lemma~\ref{lem:rate-filtered}, where $F_j$ is written $\Phi_j$), and is dominated by a polynomial in them; and $\pi_r(v_r)$ is continuous in the target. Also $\mathcal A_0(b,k,u)\le k^2H\to0$. By \eqref{eq:remainder-polar}, dominated convergence in \eqref{eq:remainder-diff} gives
\[
\nu_{\cand}^{\mathrm{near},r_0}(\ell)-c_{d,L}\ell^{-1/3}\to\int_0^{r_0}\!\!\int\!\!\int r^{d-1}\Psi_0(b,ru)\,d\sigma\,db\,dr=\int_{0<\dist(0,y)<r_0}\int_\R\Psi_0\,db\,dy,
\]
and the majorant gives the bound in \eqref{eq:remainder-Rp3}.

\emph{The far candidate density.} On the compact set $\{\dist(0,y)\ge r_0\}$, Proposition~\ref{prop:remainder-Z}(iii) bounds $\Psi^{\cand}_\ell$ by an integrable function of $b$, uniformly in $y$ and in $0<\ell\le1$, and $\Psi^{\cand}_\ell\to\Psi_0$ pointwise by Proposition~\ref{prop:remainder-Z}(ii). By \eqref{eq:remainder-far} and dominated convergence, $\nu_{\cand}^{\mathrm{far},r_0}(\ell)\to\int_{\dist(0,y)\ge r_0}\int\Psi_0$. Adding the near limit, $\nu_{\cand}(\ell)-c_{d,L}\ell^{-1/3}\to B_{d,L}$, which is finite by Proposition~\ref{prop:remainder-Z}(i) (the cutoff boundary is spatially null). This is \eqref{eq:remainder-Rp1}.

\emph{The elder density.} $\nu_{\eld}-c_{d,L}\ell^{-1/3}=(\nu_{\cand}-c_{d,L}\ell^{-1/3})-\rho_{\rej}\to B_{d,L}-B_{d,L}=0$ by \eqref{eq:remainder-Rp1} and the limit in Proposition~\ref{prop:remainder-Z}(i), which \S\ref{ssec:remainder-Z} derives from \eqref{eq:remainder-Rp1} and Theorem~\ref{thm:remainder-rate}; equivalently, \eqref{eq:remainder-Rp2} is the case $o(1)$ of \eqref{eq:remainder-Rpp}. This is \eqref{eq:remainder-Rp2}.
\end{proof}

The proof shows that the constant $C$ in \eqref{eq:remainder-Rp3} depends on $d$ and $L$ only: it comes from the majorant $H(b,\ell/r^3)\le Ce^{-cb^2}$ on $r<r_0$ and from $\int_{r_0}^\infty r^{-2}(\ell/r^3)^2H\,dr\le C\ell^2r_0^{-7}$.

\begin{remark}[a signed near bound]\label{rem:remainder-signed}
Always $\mathcal A_r\ge0$ ($\pi_r\ge0$ and $F_j\ge0$) and $\mathcal A_0\le k^2H$. Put $r_1:=\min(\ell^{1/4},r_0)$. In \eqref{eq:remainder-diff}, on $0<r<r_1$ Lemma~\ref{lem:remainder-even} gives $r^{-2}(\mathcal A_r-\mathcal A_0)\ge-Ce^{-cb^2}$, contributing at least $-Cr_1$; on $r_1\le r<r_0$, $r^{-2}(\mathcal A_r-\mathcal A_0)\ge-\ell^2r^{-8}H$, contributing at least $-C\ell^2r_1^{-7}/7=-C\ell^{1/4}/7$ when $r_1=\ell^{1/4}$ (and nothing when $r_1=r_0$); the far term is at least $-C\ell^2r_0^{-7}$. Hence
\[
\nu_{\cand}^{\mathrm{near},r_0}(\ell)-c_{d,L}\ell^{-1/3}\ge-C\bigl(\min(\ell^{1/4},r_0)+\ell^2r_0^{-7}\bigr).
\]
A lower bound $-O(\ell^2r_0^{-7})$ would need a positive lower bound on the near-diagonal equal-height kernel, which is not available here.
\end{remark}

\begin{remark}[where the bounded remainder goes]\label{rem:remainder-where}
For the candidate density the bounded remainder of Theorem~R is real: what the leading law $c_{d,L}\ell^{-1/3}$, the small-separation contact mass, misses is the equal-height mass $B_{d,L}$, spread over all separations, and by Proposition~\ref{prop:remainder-Z}(i) all of it is rejected mass. For the elder density, the three bounded contributions in the proof of Theorem~R --- the first-order contact monomial $k^{-2/3}\,rk=\ell^{1/3}$, the terms proportional to the cutoff $r_0$, and the far term --- are, respectively, an artefact of bounding the two endpoints separately (Lemma~\ref{lem:remainder-even}), the rejected mass at separations below $r_0$, and the far rejected mass; none of them is elder mass in the limit.
\end{remark}

\subsection{The rate}\label{ssec:remainder-rate}

\begin{theorem}[Theorem R$^{++}$: rate of the elder remainder]\label{thm:remainder-rate}
For every fixed $d\ge2$ and $L>0$, and for the canonical marked Kac--Rice density version of the elder density, there are $C<\infty$ and $\ell_0\in(0,1]$ depending on $d,L$ only such that
\begin{equation}\label{eq:remainder-Rpp}
\bigl|\nu_{\eld}(\ell)-c_{d,L}\,\ell^{-1/3}\bigr|\le C\ell^{1/4},\qquad 0<\ell\le\ell_0;
\end{equation}
equivalently, $\nu_{\eld}(\ell)=c_{d,L}\,\ell^{-1/3}\bigl(1+O(\ell^{7/12})\bigr)$ as $\ell\downarrow0$.
\end{theorem}

No sharpness of the exponent is asserted here (see Remark~\ref{rem:remainder-exponent}), and the theorem gives no numerical $C$ or $\ell_0$, no rate for the candidate remainder in \eqref{eq:remainder-Rp1}, no rate for $\rho_{\rej}\to B_{d,L}$, and no uniformity in $d$ or $L$. The proof uses three estimates: a deterministic barrier (Lemma~\ref{lem:remainder-barrier}), which bounds the elder density at intermediate and at fixed separations (Lemmas~\ref{lem:remainder-B} and~\ref{lem:remainder-far}), and a majorant for the near-pair weight at small gap marks (Lemma~\ref{lem:remainder-W}).

\begin{lemma}[deterministic barrier]\label{lem:remainder-barrier}
Let $f\in C^3(\T^d_L)$ and $M\in\T^d_L$ with $f(M)=b$, $\nabla f(M)=0$ and $H_M=D^2f(M)\prec0$. Put $\lambda:=\lambda_{\min}(-H_M)$, $K:=1+\|D^2f\|_{\infty,\mathrm{op}}+\|D^3f\|_{\infty,\mathrm{op}}$ (suprema over $\T^d_L$), $a_L:=\min(1,L/4)$ and $c_L:=a_L^2/3$. Then
\begin{equation}\label{eq:remainder-barrier}
b-d_f(M)\ge c_L\lambda^3/K^2 .
\end{equation}
Consequently, if $d_f(M)=b-\ell$ with $\ell>0$, then $\lambda_{\min}(-H_M)\le(\ell K^2/c_L)^{1/3}$ and $|\det H_M|\le c_L^{-1/3}\ell^{1/3}K^{d-1/3}$.
\end{lemma}

\begin{proof}
For a unit vector $v$ and $0<t\le R:=a_L\lambda/K$, Taylor's theorem gives $f(M+tv)\le b-\frac\lambda2t^2+\frac K6t^3\le b-\frac\lambda3t^2$, since $t\le\lambda/K$. The ball $B(M,R)$ is embedded ($R\le L/4$), no point of it lies above $b$, and every path from $M$ to a point above $b$ crosses its boundary sphere, on which $f\le b-\lambda R^2/3$. Hence $d_f(M)\le b-\lambda R^2/3=b-c_L\lambda^3/K^2$. For the consequence, $-H_M\succ0$ has eigenvalues $\lambda=\lambda_1\le\dots\le\lambda_d\le K$, so $|\det H_M|\le\lambda K^{d-1}\le c_L^{-1/3}\ell^{1/3}K^{d-1/3}$.
\end{proof}

\begin{lemma}[short elder bars at near separations are born at soft maxima]\label{lem:remainder-B}
For $0<r\le r_R$, $b\in\R$ and $k>0$, with $\ell=kr^3$,
\begin{equation}\label{eq:remainder-B1}
\E_Q\bigl[(W_r/r^2)\,e\bigr]\le C\ell^{1/3}r^{-1}(k+r)P^N,
\end{equation}
and consequently, for $0<\rho\le r_0\le r_R$ and $\ell>0$,
\begin{equation}\label{eq:remainder-B2}
\nu_{\eld}^{[\rho,r_0)}(\ell)=\int_\rho^{r_0}\!\!\int_\R\int_{S^{d-1}}r^{-2}\mathcal A^{\eld}_r(b,\ell/r^3,u)\,d\sigma\,db\,dr\le C\ell^{1/3}\bigl(\rho^{-1}+\ell\rho^{-5}\bigr).
\end{equation}
\end{lemma}

\begin{proof}
Under $Q$ the field is almost surely $C^\infty$. On $\{W_r>0\}\cap\{e=1\}$ we have $\nabla f(M)=0$, $f(M)=b$, $H_M\prec0$ and $d_f(M)=f(S)=b-\ell$, so Lemma~\ref{lem:remainder-barrier} gives $|\det H_M|\le c_L^{-1/3}\ell^{1/3}K^{d-1/3}$. On the support of $W_r$, $W_r/r^2=(|\det H_M|/r)\,F_{d-1}(K_S)$, and $F_{d-1}(K_S)\le C(k+r)T^d$ (Proposition~\ref{prop:remainder-R}(iv)). Hence $(W_r/r^2)e\le C\ell^{1/3}r^{-1}(k+r)K^{d-1/3}T^d$ everywhere. Since $K\le C_d(1+\|F_r\|_{C^3})$ has $\|K\|_{L^p}\le C_pP$ by \eqref{eq:remainder-R4}, H\"older's inequality gives $\E_Q[K^{d-1/3}T^d]\le CP^N$, which is \eqref{eq:remainder-B1}. The identity in \eqref{eq:remainder-B2} is \eqref{eq:remainder-sep}. By Proposition~\ref{prop:remainder-R}(ii), with $k=\ell/r^3$, $(1+|b|+k)^N\le2^N(1+|b|)^N(1+k)^N$ and $(1+k)^Ne^{-ck^2}\le C$,
\[
r^{-2}\mathcal A^{\eld}_r(b,\ell/r^3,u)\le C\ell^{1/3}r^{-3}(\ell r^{-3}+r)P^Ne^{-c(b^2+k^2)}\le C\ell^{1/3}(\ell r^{-6}+r^{-2})(1+|b|)^Ne^{-cb^2},
\]
and $\int_\rho^{r_0}(\ell r^{-6}+r^{-2})\,dr\le\ell\rho^{-5}/5+\rho^{-1}$.
\end{proof}

\begin{lemma}[far elder density, crude rate]\label{lem:remainder-far}
For every fixed $r_0\in(0,L/2]$ there is $C=C(d,L,r_0)$ with
\begin{equation}\label{eq:remainder-F}
\nu_{\eld}^{\mathrm{far},r_0}(\ell)\le C\ell^{1/3},\qquad 0<\ell\le1 .
\end{equation}
\end{lemma}

\begin{proof}
Let $D_{r_0}=\{y:\dist(0,y)\ge r_0\}$ (compact) and use the far identity \eqref{eq:remainder-far}. The covariance $\Gamma_y=\Cov(O_y)$ is positive definite for every $y\ne0$ (finite-jet nondegeneracy at distinct sites, \S\ref{sec:setting}; \cite[\S2]{parentP}; Lemma~\ref{lem:parent-jetrank}) and continuous in $y$, so $\lambda_{\min}(\Gamma_y)\ge c_0>0$ on $D_{r_0}$ \cite[\S14]{parentP} (Lemma~\ref{lem:parent-far}). Consequently:
(a) $p_y(v)\le(2\pi c_0)^{-(d+1)}\exp(-|v|^2/(2\tr\Gamma_y))\le Ce^{-cb^2}$ at $v=v_{b,\ell}$ ($|v_{b,\ell}|^2\ge b^2$, $\tr\Gamma_y\le C$);
(b) with $\psi^y:=\Gamma_y^{-1}\Cov(O_y,f(\cdot))$, $2d+2$ smooth functions with $\|\psi^y_j\|_{C^3}\le C$ uniformly on $D_{r_0}$, the field decomposes as $f=\sum_jO_y^j\psi^y_j+g_y$ with $g_y$ a centred Gaussian field independent of $O_y$, and under $Q_{y,b,\ell}$ the field has the law of $\sum_jv_{b,\ell}^j\psi^y_j+g_y$. Since $\|g_y\|_{C^3}\le\|f\|_{C^3}+C|O_y|$ and all moments of $\|f\|_{C^3}$ and $|O_y|$ are finite, $\E\|g_y\|_{C^3}^p\le C_p$ uniformly on $D_{r_0}$. For $0<\ell\le1$ ($|v_{b,\ell}|\le2|b|+1$) this gives $K\le C(1+|b|+\|g_y\|_{C^3})$ and
\begin{equation}\label{eq:remainder-Kmom}
\E_{Q_{y,b,\ell}}[K^p]\le C_p(1+|b|)^p\qquad(p\ge1,\ y\in D_{r_0},\ 0<\ell\le1).
\end{equation}
On $\{W>0\}\cap\{e_y=1\}$ the point $0$ is a nondegenerate maximum with $f(0)=b$ and $d_f(0)=f(y)=b-\ell$, so Lemma~\ref{lem:remainder-barrier} gives $|\det H_0|\le c_L^{-1/3}\ell^{1/3}K^{d-1/3}$, while $|\det H_y|\le K^d$. So $We_y\le C\ell^{1/3}K^{2d-1/3}$ and, by \eqref{eq:remainder-Kmom}, $\E_{Q_{y,b,\ell}}[We_y]\le C\ell^{1/3}(1+|b|)^{2d}$. With (a) in \eqref{eq:remainder-far}, $\nu_{\eld}^{\mathrm{far},r_0}(\ell)\le C\ell^{1/3}\operatorname{vol}(D_{r_0})\int_\R(1+|b|)^{2d}e^{-cb^2}db\le C\ell^{1/3}$.
\end{proof}

Lemma~\ref{lem:remainder-far} needs no Hessian density and suffices for this section, because the far term is not the bottleneck of the estimate below. Keeping the determinant weight through a spectral integration gives the sharper rate $\ell^{2/3}$, which Section~\ref{sec:third} uses.

\begin{proposition}[far elder density, rate $\ell^{2/3}$]\label{prop:remainder-farrate}
For every fixed $\rho\in(0,L/2]$ there is $C=C(d,L,\rho)$ with
\begin{equation}\label{eq:remainder-farrate}
\nu_{\eld}^{\mathrm{far},\rho}(\ell)\le C\ell^{2/3},\qquad0<\ell\le1 .
\end{equation}
\end{proposition}

\begin{proof}
Let $D_\rho=\{y:\dist(0,y)\ge\rho\}$ and use the far identity \eqref{eq:remainder-far}. Append to $O_y$ the independent entries of $H_0$ and $H_y$: $V_y=(O_y,H_0,H_y)$, with $2(d+1)+d(d+1)$ coordinates. These are distinct derivative evaluations at two distinct sites, so the covariance of $V_y$ is positive definite (\S\ref{sec:setting}; \cite[\S2]{parentP}; Lemma~\ref{lem:parent-jetrank}). It is continuous in $y$, so its eigenvalues lie in a fixed compact subinterval of $(0,\infty)$ on the compact set $D_\rho$. Two consequences follow.

(a) Since $|v_{b,\ell}|^2\ge b^2$, uniformly for $y\in D_\rho$, $b\in\R$ and $\ell>0$,
\begin{equation}\label{eq:remainder-farrate-dens}
p_y(v_{b,\ell})\,p_{(H_0,H_y)\mid O_y=v_{b,\ell}}(H_0,H_y)\le Ce^{-c(b^2+|H_0|^2+|H_y|^2)} .
\end{equation}

(b) Regressing the field on $V_y$, as in (b) of the proof of Lemma~\ref{lem:remainder-far}, gives $f=m+g$ with $g$ independent of $V_y$, $\E\|g\|_{C^3}^q\le C_q$ and $\|m\|_{C^3}\le C(1+|b|+|H_0|+|H_y|)$ for $0<\ell\le1$. So, with $K$ as in Lemma~\ref{lem:remainder-barrier}, for every $q\ge1$
\begin{equation}\label{eq:remainder-farrate-K}
\E\bigl[K^q\bigm|V_y=(v_{b,\ell},H_0,H_y)\bigr]\le C_q(1+|b|+|H_0|+|H_y|)^q .
\end{equation}

Fix $y\in D_\rho$, $b\in\R$ and $0<\ell\le1$. On $\{W>0\}\cap\{e_y=1\}$ the point $0$ is a nondegenerate maximum with $f(0)=b$ and $d_f(0)=f(y)=b-\ell$. So Lemma~\ref{lem:remainder-barrier}, with $\lambda=\lambda_{\min}(-H_0)$, gives $K\ge(c_L\lambda^3/\ell)^{1/2}$, and pathwise $We_y\le W\,\ind\{K\ge(c_L\lambda^3/\ell)^{1/2}\}$. The weight $W$ and $\lambda$ are functions of $(H_0,H_y)$. Conditioning on $V_y$ and applying Markov's inequality with \eqref{eq:remainder-farrate-K} bounds the conditional probability of the last event by $\min\{1,C_q(1+|b|+|H_0|+|H_y|)^q(\ell/(c_L\lambda^3))^{q/2}\}$. Multiply by the density \eqref{eq:remainder-farrate-dens} and bound $(1+|b|+|H_0|+|H_y|)^q\le(1+|b|)^q(1+|H_0|)^q(1+|H_y|)^q$. Take the two Hessian factors outside the minimum, since $\min\{1,uv\}\le v\min\{1,u\}$ for $v\ge1$, and integrate $H_y$ against $|\det H_y|(1+|H_y|)^qe^{-c|H_y|^2}$, a constant. For $H_0$ use the ordered spectral coordinates $-H_0=O\,\mathrm{diag}(\lambda_1\le\dots\le\lambda_d)\,O^{\top}$ on the cone $\{H_0\prec0\}$, whose Lebesgue Jacobian is a constant times $\prod_{i<j}(\lambda_j-\lambda_i)$, with finite angular volume. With $|\det H_0|=\prod_j\lambda_j$, bound the Vandermonde factor and $\prod_{j\ge2}\lambda_j$ by a polynomial in $|\lambda|$, extend $\lambda_2,\dots,\lambda_d$ to $(0,\infty)$ and integrate them out. This gives
\begin{equation}\label{eq:remainder-farrate-int}
p_y(v_{b,\ell})\,\E_{Q_{y,b,\ell}}[We_y]\le Ce^{-cb^2}\int_0^\infty\lambda_1\min\Bigl\{1,A\Bigl(\frac{\ell}{\lambda_1^3}\Bigr)^{q/2}\Bigr\}(1+\lambda_1)^Ne^{-c\lambda_1^2}\,d\lambda_1,\qquad A=C(1+|b|)^q .
\end{equation}
Take $q=2$, and let $u=(A\ell)^{1/3}$ be the point where the two arguments of the minimum agree. If $u\ge1$, bound the minimum by $1$: the integral is at most $C\le Cu^2$. If $u<1$, split at $u$. On $[0,u]$ the integral is at most $2^{N-1}u^2$. On $[u,\infty)$ it is at most
\[
A\ell\int_u^\infty\lambda_1^{-2}(1+\lambda_1)^Ne^{-c\lambda_1^2}\,d\lambda_1\le CA\ell\,(u^{-1}+1)=Cu^3(u^{-1}+1)\le2Cu^2 .
\]
In every case the integral in \eqref{eq:remainder-farrate-int} is at most $Cu^2=CA^{2/3}\ell^{2/3}=C(1+|b|)^{4/3}\ell^{2/3}$. Integrating over $b\in\R$ and $y\in D_\rho$ in \eqref{eq:remainder-far} proves \eqref{eq:remainder-farrate}.
\end{proof}

The factor $\lambda_1$ in \eqref{eq:remainder-farrate-int} is the soft eigenvalue of $|\det H_0|$. It turns the barrier constraint $\lambda_1\lesssim\ell^{1/3}$ into the mass $\ell^{2/3}$; without the determinant weight the same argument would give only $\ell^{1/3}$, as in Lemma~\ref{lem:remainder-far}. At $d=3$, $L=24$ and $\rho<12$, Proposition~\ref{prop:remainder-farrate} sharpens the bound of Proposition~\ref{prop:limit-far}.

\begin{lemma}[sign-window majorant for the near-pair weight]\label{lem:remainder-W}
For $0<r\le r_R$, $b\in\R$ and $k\ge0$ (the pinned family at $k=0$ included),
\begin{equation}\label{eq:remainder-W1}
\E_Q[W_r/r^2]\le C\bigl(k^2+r^4(1+k)^2\bigr)P^N .
\end{equation}
Consequently, for $0<\rho\le r_0^*$ and $0<\ell\le1$,
\begin{align}
\bigl|\nu_{\cand}^{\mathrm{near},\rho}(\ell)-c_{d,L}\,\ell^{-1/3}\bigr|&\le C\bigl(\min(\ell^{1/4},\rho)+\rho^3+\ell^2\rho^{-7}\bigr),\label{eq:remainder-W2}\\
0\le\rho_{\rej}^{\mathrm{near},\rho}(\ell)&\le C\bigl(\ell^{1/4}+\rho^3\bigr),\label{eq:remainder-W3}\\
\int_{0<\dist(0,y)<\rho}\int_\R\Psi_0(b,y)\,db\,dy&\le C\rho^3 .\label{eq:remainder-W4}
\end{align}
\end{lemma}

Since $k^2+r^4(1+k)^2\le2(1+r_R)^2(k+r)^2$, the bound \eqref{eq:remainder-W1} is never worse than the first bound in \eqref{eq:remainder-K}, and for $k<r$ it is better by the factor $k^2/r^2+r^2$. Bounds \eqref{eq:remainder-W1}, \eqref{eq:remainder-W3} and \eqref{eq:remainder-W4} use only Steps~1--3 of the proof of Lemma~\ref{lem:remainder-even}; \eqref{eq:remainder-W2} uses the full Lemma~\ref{lem:remainder-even}. In particular \eqref{eq:remainder-W4} says that the radial density $r^{d-1}\Psi_0(b,ru)$ of the equal-height kernel is $O(r^2)$ near the diagonal.

\begin{proof}
\emph{Proof of \eqref{eq:remainder-W1}.} Work under the coupling as in the proof of Lemma~\ref{lem:remainder-even}, with $A_0=D_y^2F_0(0)$ and $\Delta:=\det A_0$; nothing below divides by $k$. Steps~1--3 there hold pathwise for $0<r\le r_R$; multiplying out $\det K_i=\alpha_i\det A_i-rq_i$ with $\tau:=\tr(\adj(A_0)B)$, $\det A_i=\Delta\mp\tfrac r2\tau+E'_i$, $|E'_i|\le Cr^2T^m$ (by \eqref{eq:remainder-A}), $\alpha_i=\mp6k+\tfrac r{12}f_4+\rho_i$ (by \eqref{eq:remainder-alpha}) and $q_i=q+O(rT^{m+1})$ (by \eqref{eq:remainder-beta}) gives
\begin{equation}\label{eq:remainder-detK}
\det K_M=-6k\Delta+rY+R_M,\qquad \det K_S=+6k\Delta+rY+R_S,\qquad |R_M|+|R_S|\le Cr^2(1+k)T^{N_1},
\end{equation}
with $Y$ as in \eqref{eq:remainder-Y} and
\[
\begin{aligned}
R_{M}&=-6kE'_M-\tfrac{r^2}{24}f_4\tau+\tfrac r{12}f_4E'_M+\rho_M\det A_M-r(q_M-q),\\
R_{S}&=6kE'_S+\tfrac{r^2}{24}f_4\tau+\tfrac r{12}f_4E'_S+\rho_S\det A_S-r(q_S-q);
\end{aligned}
\]
each term is at most $Cr^2(1+k)T^{N_1}$, and the $r$-linear terms $\tfrac r{12}f_4\Delta$, $3kr\tau$ and $-rq$ are the same at both endpoints. On the support of $W_r$, $\Ind K_M=d$ and $\Ind K_S=d-1$, so $\det K_M$ has the sign $(-1)^d$, $\det K_S$ the sign $(-1)^{d-1}$, and $\det K_M\det K_S<0$. With $w:=rY$, the product $(w-6k\Delta+R_M)(w+6k\Delta+R_S)$ is negative only when $w$ lies strictly between $6k\Delta-R_M$ and $-6k\Delta-R_S$; hence
\[
|w|\le6k|\Delta|+\max(|R_M|,|R_S|),\qquad |\det K_i|\le12k|\Delta|+2\max(|R_M|,|R_S|)\quad(i=M,S),
\]
and on the support of $W_r$
\begin{equation}\label{eq:remainder-window}
W_r/r^2=|\det K_M\det K_S|\le\bigl(12k|\Delta|+2Cr^2(1+k)T^{N_1}\bigr)^2\le C\bigl(k^2\Delta^2+r^4(1+k)^2T^{2N_1}\bigr).
\end{equation}
Completing the square gives the same bound more directly: $\det K_M\det K_S=\bigl(w+\tfrac12(R_M+R_S)\bigr)^2-\bigl(6k\Delta+\tfrac12(R_S-R_M)\bigr)^2$, so on the support of $W_r$ we have $W_r/r^2\le\bigl(6k|\Delta|+\tfrac12(|R_M|+|R_S|)\bigr)^2$. Off the support $W_r=0$. Since $|\Delta|\le\|A_0\|^m\le T^m$, $\E\Delta^2\le CP^{2m}$ and $\E T^{2N_1}\le CP^{2N_1}$; $W_r$ is a functional of $F_r$, whose law is $Q$, so the pathwise bound on the coupling space bounds $\E_Q$. This is \eqref{eq:remainder-W1}.

\emph{Proof of \eqref{eq:remainder-W2}.} Put $r_1:=\ell^{1/4}$, so that $r\le r_1$ if and only if $k=\ell/r^3\ge r$, and use \eqref{eq:remainder-diff} with $r_0=\rho$. On $0<r<\min(r_1,\rho)$, Lemma~\ref{lem:remainder-even} gives $|r^{-2}(\mathcal A_r-\mathcal A_0)|\le H(b,k)\le Ce^{-cb^2}$, contributing at most $C\min(r_1,\rho)$. On $r_1\le r<\rho$ (empty if $\rho\le r_1$), $|\mathcal A_r-\mathcal A_0|\le\mathcal A_r+\mathcal A_0$ with $\mathcal A_r\le C(k^2+r^4(1+k)^2)P^Ne^{-c(b^2+k^2)}$ by \eqref{eq:remainder-W1} and Proposition~\ref{prop:remainder-R}(ii), and $\mathcal A_0\le k^2H$; so $r^{-2}|\mathcal A_r-\mathcal A_0|\le C(\ell^2r^{-8}+r^2)(1+|b|)^Ne^{-cb^2}$, which integrates to at most $C(\ell^2r_1^{-7}/7+\rho^3/3)=C(\ell^{1/4}/7+\rho^3/3)$. The tail term is at most $C\ell^2\rho^{-7}$ as in the proof of Theorem~\ref{thm:remainder-vanish}.

\emph{Proof of \eqref{eq:remainder-W3}.} By \eqref{eq:remainder-sep}, $\rho_{\rej}^{\mathrm{near},\rho}=\int_0^\rho\int\int r^{-2}\mathcal A^{\rej}_r$. On $0<r\le\min(r_1,\rho)$, Proposition~\ref{prop:remainder-K} (valid as $r\le k$ and $r\le\rho\le r_0^*$) gives $r^{-2}\mathcal A^{\rej}_r\le Cr^{-2}(r^3/k)P^Ne^{-c(b^2+k^2)}=C(r^4/\ell)P^Ne^{-c(b^2+k^2)}$, which integrates to at most $Cr_1^5/(5\ell)=C\ell^{1/4}/5$. On $r_1<r\le\rho$, $1-e\le1$ and \eqref{eq:remainder-W1}, as in the proof of \eqref{eq:remainder-W2}, give at most $C(\ell^{1/4}+\rho^3)$.

\emph{Proof of \eqref{eq:remainder-W4}.} By \eqref{eq:remainder-polar}, $r^{d-1}\Psi_0(b,ru)=r^{-2}\mathcal A_r(b,0,u)$, and \eqref{eq:remainder-W1} at $k=0$ with Proposition~\ref{prop:remainder-R}(ii) gives $\mathcal A_r(b,0,u)\le Cr^4(1+|b|)^Ne^{-cb^2}$. So $\int_{0<\dist<\rho}\int\Psi_0=\int_0^\rho\int\int r^{-2}\mathcal A_r(b,0,u)\le C\int_0^\rho r^2dr=C\rho^3/3$.
\end{proof}

The window $|rY|\le6k|\Delta|+O(r^2)$ is where the typed weight lives at small $k$: for $k=0$ both scaled endpoint determinants are $rY+O(r^2)$ with the same $rY$, so a maximum/saddle pair can be typed only when $|Y|=O(r)$. Bound \eqref{eq:remainder-W2} needs Lemma~\ref{lem:remainder-even} on $r<\ell^{1/4}$: \eqref{eq:remainder-W1} alone gives $\int_0r^{-2}k^2\,dr=\infty$, and \eqref{eq:remainder-R11} alone gives only $O(1)$.

\begin{proof}[Proof of Theorem~\ref{thm:remainder-rate}]
Fix $0<\rho\le r_0^*$ and $0<\ell\le1$. Splitting the elder density by the separation of the pair into $(0,\rho)$, $[\rho,r_0^*)$ and $[r_0^*,\operatorname{diam}\T^d_L]$ (the boundaries are null), and using $\nu_{\eld}^{\mathrm{near},\rho}=\nu_{\cand}^{\mathrm{near},\rho}-\rho_{\rej}^{\mathrm{near},\rho}$,
\begin{equation}\label{eq:remainder-split}
\nu_{\eld}(\ell)-c_{d,L}\ell^{-1/3}=\bigl[\nu_{\cand}^{\mathrm{near},\rho}(\ell)-c_{d,L}\ell^{-1/3}\bigr]-\rho_{\rej}^{\mathrm{near},\rho}(\ell)+\nu_{\eld}^{[\rho,r_0^*)}(\ell)+\nu_{\eld}^{\mathrm{far},r_0^*}(\ell).
\end{equation}
The four terms are bounded as follows:
(i) $|\nu_{\cand}^{\mathrm{near},\rho}-c_{d,L}\ell^{-1/3}|\le C(\min(\ell^{1/4},\rho)+\rho^3+\ell^2\rho^{-7})$ by \eqref{eq:remainder-W2};
(ii) $0\le\rho_{\rej}^{\mathrm{near},\rho}\le C(\ell^{1/4}+\rho^3)$ by \eqref{eq:remainder-W3};
(iii) $0\le\nu_{\eld}^{[\rho,r_0^*)}\le C\ell^{1/3}(\rho^{-1}+\ell\rho^{-5})$ by \eqref{eq:remainder-B2} with $r_0=r_0^*\le r_R$;
(iv) $0\le\nu_{\eld}^{\mathrm{far},r_0^*}\le C\ell^{1/3}$ by \eqref{eq:remainder-F} with $r_0=r_0^*$.
Choose $\rho=\ell^{1/12}$ and $\ell_0:=\min\{1,(r_0^*)^{12}\}$, so that $\rho\le r_0^*$ for $\ell\le\ell_0$. Then $\min(\ell^{1/4},\rho)=\ell^{1/4}$, $\rho^3=\ell^{1/4}$, $\ell^2\rho^{-7}=\ell^{17/12}$, $\ell^{1/3}\rho^{-1}=\ell^{1/4}$ and $\ell^{4/3}\rho^{-5}=\ell^{11/12}$, and \eqref{eq:remainder-split} with (i)--(iv) (eight terms, five of them equal to $\ell^{1/4}$) gives
\[
|\nu_{\eld}(\ell)-c_{d,L}\ell^{-1/3}|\le C\bigl(5\ell^{1/4}+\ell^{17/12}+\ell^{11/12}+\ell^{1/3}\bigr)\le8C\ell^{1/4},\qquad0<\ell\le\ell_0,
\]
which is \eqref{eq:remainder-Rpp}. Dividing by $c_{d,L}\ell^{-1/3}$ ($c_{d,L}>0$) gives the relative form.
\end{proof}

\begin{remark}[what fixes the exponent]\label{rem:remainder-exponent}
In \eqref{eq:remainder-split} the leading $\rho$-dependence is $\rho^3$ (the candidate excess and the rejected mass at separations in $(\ell^{1/4},\rho)$, both $O(r^2)$ per unit separation by Lemma~\ref{lem:remainder-W}) against $\ell^{1/3}\rho^{-1}$ (Lemma~\ref{lem:remainder-B} above $\rho$); they balance at $\rho=\ell^{1/12}$, where both equal $\ell^{1/4}$. But $\ell^{1/4}$ also appears three times independently of $\rho$, all from the transition $r=k$ at separation $\ell^{1/4}$, below which the pinned pair is a cap pair: the excess $C\min(\ell^{1/4},\rho)$ from Lemma~\ref{lem:remainder-even}, the loss from Proposition~\ref{prop:remainder-K} at $r\le k$, and the tail $\int_{r>\ell^{1/4}}r^{-2}k^2\,dr$ of the leading law, $\mathcal A_0\asymp k^2$. The order $\ell^{1/4}$ is attained, with a negative coefficient, by Theorem~\ref{thm:cusp-CU}: at the cusp scale $r\asymp k\asymp\ell^{1/4}$ the contact law overcounts. Note also that \eqref{eq:remainder-Rpp} does not use Proposition~\ref{prop:remainder-Z}(i); together with \eqref{eq:remainder-Rp1} it re-derives $\rho_{\rej}(\ell)\to B_{d,L}$.
\end{remark}

\subsection{The cap-region bound}\label{ssec:remainder-capK}

This subsection proves the second bound of Proposition~\ref{prop:remainder-K}, which is also part of Theorem~K of record [R7] of the main text. The proof is that of Lemma~\ref{lem:rate-caploss}, with two changes: the axial entries of the endpoint Hessians are bounded by \eqref{eq:rate-52} instead of \eqref{eq:rate-51}, which keeps a factor $k^2$ in the weight, and the pin density is left out, so that the density of the transverse Hessian is a Gaussian with a shifted centre.

\begin{lemma}[the cap-region bound]\label{lem:remainder-capK}
Let $r_0$ be as in \eqref{eq:rate-r0}. There are $C$ and $N$ such that, for $0<r\le r_0$, $b\in\R$, $k\ge r$ and every frame,
\begin{equation}\label{eq:remainder-capK}
\E_Q\Bigl[\frac{W_r}{r^2}\,\ind_{G_r^c}\Bigr]\le C\,\frac{r^3}k\,P^N .
\end{equation}
\end{lemma}

\begin{proof}
Put $\delta=r/k\le1$. The law $Q$ is the law $\E^0$ of Section~\ref{sec:rate}, $G_r$ is the cap event \eqref{eq:rate-G}, and $\alpha_i,\beta_i,A_i$ are as in Lemma~\ref{lem:rate-endpoint}.

\emph{The weight with the axial entries at their true size.} Put $a:=6k+rM_4/2$, so that $|\alpha_M|,|\alpha_S|\le a$ by \eqref{eq:rate-52}. By \eqref{eq:rate-53}, $|\det H_i|/r\le|\alpha_i|\,|\det A_i|+r\|\beta_i\|^2\|\adj(A_i)\|$, and \eqref{eq:rate-51} gives
\begin{equation}\label{eq:remainder-capW}
\frac{W_r}{r^2}\le\Bigl(a\,\|f\|_{C^2}^m+\frac r4\,M_3^2\,\|f\|_{C^2}^{m-1}\Bigr)^2\le C(k+r)^2\bigl(1+\|f\|_{C^4}\bigr)^{2d+2}.
\end{equation}
On the typed support $H_M\prec0$, so $\alpha_M<0$ and $B:=-A_M\succ0$; let $0<\lambda_1\le\dots\le\lambda_m=\Lambda$ be the eigenvalues of $B$. By \eqref{eq:rate-53}, $|\det H_M|=r\bigl(|\alpha_M|\det B-r\beta_M^{\top}\adj(B)\beta_M\bigr)\le ra\prod_j\lambda_j$, as $\adj(B)\succeq0$. By \eqref{eq:rate-51} and Weyl's inequality the eigenvalues of $-A_S$ lie within $rM_3$ of those of $B$, so $|\det A_S|\le\prod_j(\lambda_j+rM_3)$ and $\|\adj(A_S)\|\le\prod_{j\ge2}(\lambda_j+rM_3)$; with $\|\beta_S\|\le M_3/2$, \eqref{eq:rate-53} gives $|\det H_S|\le r\bigl[a(\lambda_1+rM_3)+\frac r4M_3^2\bigr]\prod_{j\ge2}(\lambda_j+rM_3)$. Hence
\begin{equation}\label{eq:remainder-capenv}
\frac{W_r}{r^2}\le a\,\lambda_1\Bigl[a\lambda_1+a\,rM_3+\frac r4\,M_3^2\Bigr]\prod_{j=2}^m\lambda_j\,(\lambda_j+rM_3);
\end{equation}
for $a=M_3/2$ this is \eqref{eq:rate-envelope}.

\emph{Residual and the density of $A_M$.} As in the proof of Lemma~\ref{lem:rate-caploss}, Lemma~\ref{lem:parent-residual} writes $f$ under $Q$ as a regression mean, linear in $(v_r,A_M)$, plus a centred residual $g$ that is independent of $A_M$ and whose $C^4$ norm has moments of every order, bounded uniformly in $r$ and in the frame. With $J=1+\|g\|_{C^4}$ and $T=P+J$ (in this proof $T$ denotes $P+J$, as in the proof of Lemma~\ref{lem:rate-caploss}, not the variable of Convention~\ref{conv:remainder-notation}(h)), on the typed support $M_3,M_4\le K_*(T+\Lambda)$ for a constant $K_*$. Under $Q$ the independent entries of $A_M$ are Gaussian. Their covariance is the Schur complement of $\Sigma_r$ in the extended covariance of Lemma~\ref{lem:rate-frame}(b), so its eigenvalues lie in a fixed compact subinterval of $(0,\infty)$. Their mean $m_A$ is linear in $v_r$, so $\|m_A\|_F\le CP$ by Lemma~\ref{lem:rate-frame}(c). Hence their density is at most $C\exp(-c\|A-m_A\|_F^2)$. In the eigenvalue coordinates of Lemma~\ref{lem:parent-eigen}, with $-A=O\operatorname{diag}(\lambda)O^{\top}$ and $O$ fixed, $\|A-m_A\|_F^2\ge\sum_i(\lambda_i+q_i)^2$, where $q_i=(O^{\top}m_AO)_{ii}$ and $|q_i|\le CP$. Since $\int_0^\infty\lambda^pe^{-c(\lambda+q)^2}d\lambda\le C_p(1+|q|)^p$, a polynomial in $\lambda_2,\dots,\lambda_m$ integrated against the Gaussian bound in these variables is at most a polynomial in $P$, uniformly in $O$; the integral over $O$ is taken last. This replaces \eqref{eq:rate-R13}, which carries the pin density.

\emph{Depth failure, $m\ge2$.} On the typed support, $G_r^c$ is the union of the depth failure $\{\lambda_1\le\frac4{3k}rM_3^2\}$ and the exception $\{rM_4>3k/10\}$. Put $\Upsilon=T+\Lambda\ge1$, which does not involve $\lambda_1$, and $D_*=\frac43K_*^2$. Depth failure gives $\lambda_1\le D_*\delta\Upsilon^2$, and as $r\le k$, $a\le\bar a:=6k+\frac12rK_*\Upsilon\le C_*k\Upsilon$; since $\bar a$ and $\Upsilon$ do not involve $\lambda_1$, we replace $a$ by $\bar a$ and $M_3$ by $K_*\Upsilon$ in \eqref{eq:remainder-capenv} before integrating in $\lambda_1$. In \eqref{eq:remainder-capenv} the product over $j\ge2$ is at most $C\Upsilon^{2m-2}$, the Vandermonde factor is at most $\Lambda^{m(m-1)/2}$, and the Gaussian factor in $\lambda_1$ is dropped. Then
\[
\bar a\int_0^{D_*\delta\Upsilon^2}\lambda\Bigl[\bar a\lambda+\bar arK_*\Upsilon+\frac r4K_*^2\Upsilon^2\Bigr]d\lambda=\frac{\bar a^2D_*^3}3\,\delta^3\Upsilon^6+\frac{\bar a^2rK_*D_*^2}2\,\delta^2\Upsilon^5+\frac{\bar arK_*^2D_*^2}8\,\delta^2\Upsilon^6 .
\]
By $\bar a\le C_*k\Upsilon$, $\delta=r/k$ and $r^3\le P\,r^3/k$, each term is at most $C(r^3/k)P\,\Upsilon^8$. The integration in $\lambda_2,\dots,\lambda_m$ and over $O$, and the moments of $J$, give
\[
\E_Q\Bigl[\frac{W_r}{r^2}\,\ind_{\{\lambda_1\le\frac4{3k}rM_3^2\}}\Bigr]\le C\,\frac{r^3}k\,P^N .
\]

\emph{Depth failure, $m=1$.} Now $\lambda_1=\Lambda=\lambda$, and by \eqref{eq:rate-R15} depth failure gives $\lambda\le4D_*\delta T^2$ or $\lambda>1/(4D_*\delta)$. On the first branch $M_3,M_4\le K_*(T+4D_*T^2)\le CT^2$ and $a\le CkT^2$, and integrating \eqref{eq:remainder-capenv} over $[0,4D_*\delta T^2]$ as above gives at most $C(r^3/k)P\,T^{10}$, hence $C(r^3/k)P^N$ after the moments of $J$. On the second branch $1\le(4D_*\delta\lambda)^4$, so by \eqref{eq:remainder-capW}, $\lambda\le\|f\|_{C^2}$ and \eqref{eq:rate-R4} (under $Q$ the field is $F_r$ of Lemma~\ref{lem:rate-coupling}, for every $0<r\le r_0$),
\[
\E_Q\Bigl[\frac{W_r}{r^2}\,\ind_{\{\lambda>1/(4D_*\delta)\}}\Bigr]\le(4D_*\delta)^4\,\E_Q\Bigl[\frac{W_r}{r^2}\,\lambda^4\Bigr]\le C\delta^4(k+r)^2P^N\le C\,\frac{r^3}k\,P^N,
\]
since $\delta^4(k+r)^2\le4\delta^4k^2=4\delta\,r^3/k$.

\emph{Exception.} On $\{rM_4>3k/10\}$, $1\le(10\delta M_4/3)^4$, so by \eqref{eq:remainder-capW} and \eqref{eq:rate-R4},
\[
\E_Q\Bigl[\frac{W_r}{r^2}\,\ind_{\{rM_4>3k/10\}}\Bigr]\le\Bigl(\frac{10\delta}3\Bigr)^4\E_Q\Bigl[\frac{W_r}{r^2}\,M_4^4\Bigr]\le C\delta^4(k+r)^2P^N\le C\,\frac{r^3}k\,P^N .
\]
Adding the bounds for depth failure and for the exception proves \eqref{eq:remainder-capK}.
\end{proof}

With the axial bound \eqref{eq:rate-51} in place of \eqref{eq:rate-52}, the same argument gives only $C(r/k)^3P^N$, the analogue without the pin density of the bound \eqref{eq:rate-R12} of Lemma~\ref{lem:rate-caploss}; the factor $k^2$ kept here is the factor $(6k)^2$ of $z_0$ in \eqref{eq:rate-z0}, and it is what the uses of Proposition~\ref{prop:remainder-K} in Lemma~\ref{lem:remainder-W} (hence Theorem~\ref{thm:remainder-rate}), Section~\ref{sec:cusp} and \eqref{eq:third-K2} of Section~\ref{sec:third} need.

\subsection{The equal-height kernel}\label{ssec:remainder-Z}

This subsection proves items (i) and (ii) of Proposition~\ref{prop:remainder-Z}, which are also Theorem~Z of record [R7] of the main text. Item (ii) and the bounds $0<B_{d,L}<\infty$ are direct. The limit $\rho_{\rej}(\ell)\to B_{d,L}$ follows from \eqref{eq:remainder-Rp1} and \eqref{eq:remainder-Rpp}; the proof of \eqref{eq:remainder-Rp1} uses Proposition~\ref{prop:remainder-Z} only through (ii), (iii) and the finiteness of $B_{d,L}$, and the proof of \eqref{eq:remainder-Rpp} does not use it at all (cf.\ Remark~\ref{rem:remainder-exponent}), so there is no circularity.

\begin{proof}[Proof of Proposition~\ref{prop:remainder-Z}(ii)]
Fix $y\ne0$ and $b$, and let $\Gamma_y=\Cov(O_y)$ and $\mathcal C_y(z)=\Cov(f(z),O_y)$. By Lemma~\ref{lem:parent-jetrank}, applied to the values and gradients at the two distinct sites $0$ and $y$, $\Gamma_y$ is positive definite. Let $F$ be a smooth copy of the unconditioned field and put $f_\ell:=F+\mathcal C_y\Gamma_y^{-1}(v_{b,\ell}-O_y(F))$. The field $F-\mathcal C_y\Gamma_y^{-1}O_y(F)$ is uncorrelated with, hence independent of, the Gaussian vector $O_y(F)$, so $f_\ell$ has the regression law $Q_{y,b,\ell}$, and its covariance, that of $F-\mathcal C_y\Gamma_y^{-1}O_y(F)$, does not depend on $\ell$. Since $v_{b,\ell}-v_{b,0}=-\ell\,\mathsf e$, where $\mathsf e$ is the unit vector of the coordinate $f(y)$ of $O_y$,
\[
f_\ell=f_0+\ell\,h_y,\qquad h_y:=-\mathcal C_y\Gamma_y^{-1}\mathsf e,
\]
a deterministic function, smooth because each component of $\mathcal C_y$ is a derivative of $K_L$; and $O_y(h_y)=-\Gamma_y\Gamma_y^{-1}\mathsf e=-\mathsf e$, that is, $h_y(0)=0$, $\nabla h_y(0)=0$, $h_y(y)=-1$ and $\nabla h_y(y)=0$. By Lemma~\ref{lem:parent-jetrank} again, applied to the list consisting of $O_y$ and the independent entries of $H_0=D^2f(0)$ and $H_y=D^2f(y)$, the covariance of that list is positive definite, so the conditional covariance of the Hessian entries given $O_y$, its Schur complement, is positive definite too; hence under $Q_{y,b,\ell}$ the pair $(H_0,H_y)$ has a Lebesgue density on $\mathrm{Sym}(d)^2$, and the algebraic set $\{\det H_0\det H_y=0\}$, which is proper, is null. Finally, $W(f_\ell)=F_d(H_0(f_\ell))\,F_{d-1}(H_y(f_\ell))$ with $H_i(f_\ell)=H_i(f_0)+\ell D^2h_y(i)$, and $F_d$ and $F_{d-1}$ are locally Lipschitz (Lemma~\ref{lem:rate-filtered}, where $F_j$ is written $\Phi_j$); so $W(f_\ell)\to W(f_0)$ pathwise as $\ell\downarrow0$, and $W(f_\ell)\le C\bigl(1+\|f_0\|_{C^2}+\|h_y\|_{C^2}\bigr)^{2d}$ for $\ell\le1$, which has finite expectation by Lemma~\ref{lem:rate-spectrum}(a), as $\|f_0\|_{C^2}\le\|F\|_{C^2}+\|\mathcal C_y\Gamma_y^{-1}\|_{C^2}\bigl(|v_{b,0}|+|O_y(F)|\bigr)$ and $|O_y(F)|\le C\|F\|_{C^1}$. Dominated convergence gives $\E_{Q_{y,b,\ell}}[W]\to\E_{Q_{y,b,0}}[W]$, and $p_y(v_{b,\ell})\to p_y(v_{b,0})$ as $p_y$ is continuous. Hence $\Psi^{\cand}_\ell(b,y)\to\Psi_0(b,y)$.
\end{proof}

\begin{proof}[Proof of Proposition~\ref{prop:remainder-Z}(i)]
\emph{Finiteness.} Split $B_{d,L}$ at the torus distance $r_0^*$. By \eqref{eq:remainder-W4} at $\rho=r_0^*$, the part over $0<\dist(0,y)<r_0^*$ is at most $C(r_0^*)^3$; the proof of \eqref{eq:remainder-W4} uses only \eqref{eq:remainder-W1} at $k=0$, Proposition~\ref{prop:remainder-R}(ii) and \eqref{eq:remainder-polar}. On the compact set $\{\dist(0,y)\ge r_0^*\}$, item (ii) and item (iii) give $\Psi_0(b,y)=\lim_{\ell\downarrow0}\Psi^{\cand}_\ell(b,y)\le C(1+|b|)^{2d}e^{-cb^2}$, which is integrable in $b$; so that part is finite as well.

\emph{Positivity.} Fix $y\ne0$ and $b$. The density $p_y$ is positive everywhere, as $\Gamma_y$ is positive definite. Under $Q_{y,b,0}$ the pair $(H_0,H_y)$ has a positive Lebesgue density on $\mathrm{Sym}(d)^2$ (the proof of (ii)), and the set $\{H_0\prec0,\ \Ind H_y=d-1,\ \det H_y\ne0\}$, the typed set less the null set $\{\det H_0\det H_y=0\}$, is open (an index changes only where an eigenvalue vanishes) and nonempty; on it $W>0$. Hence $\E_{Q_{y,b,0}}[W]>0$ and $\Psi_0(b,y)>0$ for every $y\ne0$ and $b\in\R$, so $B_{d,L}>0$.

\emph{The limit.} By \eqref{eq:remainder-Rp1}, whose proof uses only (ii), (iii) and the finiteness of $B_{d,L}$ just proved, $\nu_{\cand}(\ell)-c_{d,L}\ell^{-1/3}\to B_{d,L}$; by \eqref{eq:remainder-Rpp}, $\nu_{\eld}(\ell)-c_{d,L}\ell^{-1/3}=O(\ell^{1/4})\to0$. Subtracting, $\rho_{\rej}(\ell)=\nu_{\cand}(\ell)-\nu_{\eld}(\ell)\to B_{d,L}$.
\end{proof}

With Lemma~\ref{lem:remainder-capK} and this subsection, every statement of Part~II rests on Sections~\ref{sec:setting}--\ref{sec:rate}, Part~III and the computer-assisted enclosures described in \S\ref{sec:cusp} and \S\ref{sec:ctwo}; no analytic statement is imported from a repository record.
